% basic nastaveni
\documentclass[12pt]{article}
\setlength{\parindent}{0pt}
\setlength{\parskip}{1em}
\usepackage[utf8]{inputenc}
\usepackage[czech]{babel}
\usepackage[T1]{fontenc}
\usepackage{graphics}
\usepackage{caption}
\usepackage{hyperref}
\usepackage{dsfont}
\usepackage{color}
\usepackage{float}
\usepackage{enumitem}
\usepackage{graphicx}

% matematicke package
\usepackage{tikz-cd}
\usetikzlibrary{babel}
\usepackage{amsmath}
\usepackage{amssymb}
\usepackage{amsfonts}
\usepackage{amssymb}
\usepackage{geometry}
\geometry{margin=2.5cm}
\usepackage{pgfplots}
\pgfplotsset{compat=1.18}
\usepackage{mathrsfs}
\usetikzlibrary{shapes.misc}
\usetikzlibrary{angles,quotes}

\title{Maturitní otázka číslo 16: Polohové a metrické úlohy ve stereometrii }
\author{}
\date{}

\begin{document}
\maketitle

\section*{Základní pojmy}
\begin{itemize} [itemsep = 2pt]
    \item \textbf{Hrana} je úsečka, ve které se protínají dvě stěny, rozdělujeme na \textbf{podstavné hrany} (tvoří obvod podstavy) a \textbf{boční hrany} (spojují podstavy, nebo tvoří "rohy" stěn)
    \item \textbf{Boční stěna} je jakákoliv stěna, která není podstavou.
    \item \textbf{Plášť} je sjednocení všech bočních stěn
    \item \textbf{Podstava} je spodní a u hranolů i horní stěna tělesa
    \item \textbf{Vrchol} je bod, ve kterém se stýkejí hrany
    \item \textbf{Kolineární body} jsou body ležící v jedné přímce
    \item \textbf{Příčka mimoběžek} je úsečka, jejíž krajní body leží na daných mimoběžkách, \textbf{Osa mimoběžek} je jediná příčka, která je k oběma mimoběžkám kolmá, její délka je zároveň nejkratší vzdáleností
    \item \textbf{Svazek rovin} je množina všech rovin, které procházejí jednou společnou přímkou (osou svazku)
    \item \textbf{Trs rovin} je množina všech rovin, které procházejí jedním společným bodem (středem trsu)
\end{itemize}

\subsection*{Úhlopříčky}
\begin{itemize} [itemsep = 2pt]
    \item \textbf{Stěnová úhlopříčka} je spojnice dvou nesousedních vrcholů ležících ve stejné stěně.
    \item \textbf{Tělesová úhlopříčka} je spojnice dvou vrcholů, které neleží ve stejné stěně. Prochází vnitřkem tělesa.
\end{itemize}

\section*{Vzájemná poloha bodů, přmek a rovin}
Přímky $p, q$
\begin{itemize} [itemsep = 2pt]
    \item \textbf{Rovnoběžné ($p \| q$)}: mohou být buď různé (žádný společný bod, kolienární vektory), nebo totožné (nekonečně společných bodů)
    \item \textbf{Různoběžné}: právě jeden společný bod, leží v jedné rovině
    \item \textbf{Mimoběžné}: Žádný společný bod, neleží v jedné rovině a směrové vektory nejsou kolineární
\end{itemize}

Přímka $p$ a rovina $\rho$
\begin{itemize} [itemsep = 2pt]
    \item \textbf{Leží v rovině ($p \in \rho)$}: Každý bod přímky je bodem roviny
    \item \textbf{Rovnoběžná ($p \| \rho$)}: žádný společný bod
    \item \textbf{Různoběžná}: právě jeden společný bod. Speciálním případem je kolmost $p \perp \rho$.
\end{itemize}

Dvě roviny $\rho , \sigma$
\begin{itemize} [itemsep = 2pt]
    \item  \textbf{Rovnoběžné}: buď různé (žádný společný bod), nebo totožné (nekonečně mnoho společných bodů)
    \item \textbf{Různoběžné:} Společné body tvoří přímku (průsečnici) 
\end{itemize}

\section*{Kritéria rovnoběžnosti}
\begin{itemize} [itemsep = 2pt]
    \item \textbf{Rovnoběžnost přímky a roviny}: Přímka $p$ je rovnoběžná s rovinou $\rho$ právě tehdy, když je rovnoběžná s alespoň jednou přímkou $q$, která v rovině $\rho$ leží.
    \item \textbf{Rovnoběžnost dvou rovin}: Rovina $\rho$ je rovnoběžná s rovinou $\sigma$, pokud $\rho$ obsahuje dvě různé různoběžky $a, b$, které jsou obě rovnoběžné s $\sigma$
\end{itemize}
\section*{Průsečík s rovinou a průsečnice dvou rovin}
Při hledání průsečnice dvou rovin hledáme dva společné body těchto rovin

U průsečíku přímky s rovinou (například v tělese) přímku proložíme pomocnou rovinou, najdeme průsečnici této pomocné roviny s cílovou rovinou a tam, kde průsečnice protne původní přímku, leží hledaný bod.
Ukážeme si to na příkladu: Najděte průsečíky přímky $PQ$ s krychlí.



\begin{tikzpicture}[
    x={(1cm,0cm)}, 
    y={(0.5cm,0.5cm)}, 
    z={(0cm,1cm)}, 
    scale=1]

    \coordinate (A) at (0,0,0);
    \coordinate (B) at (2,0,0);
    \coordinate (C) at (2,2,0);
    \coordinate (D) at (0,2,0);
    \coordinate (E) at (0,0,2);
    \coordinate (F) at (2,0,2);
    \coordinate (G) at (2,2,2);
    \coordinate (H) at (0,2,2);
    \coordinate (P) at (2,2,3);
    \coordinate (Q) at (0,0,-1);

    \draw[dashed, thick] (A) -- (D);
    \draw[dashed, thick] (D) -- (C);
    \draw[dashed, thick] (D) -- (H);

    \draw[thick] (A) -- (B) -- (C) -- (G) -- (F) -- (E) -- cycle;
    \draw[thick] (E) -- (H) -- (G);
    \draw[thick] (A) -- (E);
    \draw[thick] (B) -- (F);
    \draw[thick] (G) -- (2,2,3.5);
    \draw[thick] (A) -- (0,0,-1.5);
    \draw[blue, thick] (P) -- (Q);

    \filldraw[fill=red!30, draw=red, thick, opacity=0.5] (E) -- (G) -- (C) -- (A) -- cycle;

    \node[cross out, draw=black, thick, minimum size=6pt, inner sep=0pt, label=right:{$P$}] at (P) {};
    \node[below left] at (A) {A};
    \node[below right] at (B) {B};
    \node[right] at (C) {C};
    \node[above left] at (D) {D};
    \node[left] at (E) {E};
    \node[right] at (F) {F};
    \node[above right] at (G) {G};
    \node[above left] at (H) {H};
    \node[cross out, draw=black, thick, minimum size=6pt, inner sep=0pt, label=left:{$Q$}] at (Q) {};
    

\end{tikzpicture}

V příkladu je červeně naznačená pomocná rovina a řešením bude bod, kde se protíná úsečka $PQ$ a úsečka $EG$ a druhým průsečíkem s krychlí bude průsečík $AC$ a $PQ$. To jsou body, kde přímka $PQ$ "vchází" a "vychází" z krychle.

\section*{Řezy těles}
\textbf{Pravidla pro řezy:}
\begin{itemize} [itemsep = 2pt]
    \item Leží-li dva body v jedné stěně, spojíme je úsečkou
    \item Jsou-li dvě stěny rovnoběžné, jsou řezy v nich (průsečnice) rovnoběžné.
    \item Trsy rovin: Vezměme tři roviny: Rovinu řezu $\rho$, rovinu přední stěny $\alpha$ a rovinu boční stěny $\beta$ (rovina řezu nesmí být rovnoběžná s rovinami krychle).
    
    Víme o nich, že roviny stěn $alpha$ a $\beta$ se protínají v jedné přímce (hrana krychle, označme ji $p$) a že rovina řezu $\rho$ protíná stěnu $\alpha$ v nějaké přímce (čára řezu, označme ji $h$).

    Pokud protáhnu čáru řezu $h$ a hranu krychle $p$, tak se někde protnou, protože leží v jedné rovině a nejsou rovnoběžné. Vezměme tedy bod $X$, který leží na čáře řezu $h$, patří tedy do roviny řezu $\rho$, bod $X$ zároveň leží na přímce hrany $h$, tudíž patří i do roviny boční stěny $\beta$, mohu tedy propojit bod $X$ s dalším bodem pratřícím do řezu v rovině $\beta$.
    
    Příklad: narýsujte rovinu řezu procházející body PQF

    \begin{tikzpicture}[
    x={(1cm,0cm)}, 
    y={(0.5cm,0.5cm)}, 
    z={(0cm,1cm)}, 
    scale=1] 

    \coordinate (A) at (0,0,0);
    \coordinate (B) at (4,0,0);
    \coordinate (C) at (4,4,0);
    \coordinate (D) at (0,4,0);
    \coordinate (E) at (0,0,4);
    \coordinate (F) at (4,0,4);
    \coordinate (G) at (4,4,4);
    \coordinate (H) at (0,4,4);
    \coordinate (Q) at (4,4,2);
    \coordinate (P) at (2,0,0);
    \coordinate (X) at (4,8,0);

    \draw[dashed, thick] (A) -- (D);
    \draw[dashed, thick] (D) -- (C);
    \draw[dashed, thick] (D) -- (H);

    \draw[thick] (A) -- (B) -- (C) -- (G) -- (F) -- (E) -- cycle;
    \draw[thick] (E) -- (H) -- (G);
    \draw[thick] (A) -- (E);
    \draw[thick] (B) -- (F);
    \draw[thin] (Q) -- (F);
    \draw[thin, dashed] (Q) -- (4,10,-1);
    \draw[dashed, thin] (C) -- (4,9,0);
    \draw[dashed, thin] (X) -- (P);

    \node[cross out, draw=black, thick, minimum size=6pt, inner sep=0pt, label=right:{$P$}] at (P) {};
    \node[below left] at (A) {A};
    \node[below right] at (B) {B};
    \node[right] at (C) {C};
    \node[above left] at (D) {D};
    \node[left] at (E) {E};
    \node[above] at (F) {F};
    \node[above right] at (G) {G};
    \node[above left] at (H) {H};
    \node[below] at (X) {X};
    \node[cross out, draw=black, thick, minimum size=6pt, inner sep=0pt, label=above right:{$Q$}] at (Q) {};
    

\end{tikzpicture}

Nejdříve jsem propojil body $F, Q$, protože leží na stejné straně. Poté už je 
nutné použít trs rovin. Rovina $BFC$ a $ABC$ mají společnou průsečnici $BC$ a 
v rovině $BFC$ leží i čára hledaného řezu. Proto se po prodloužení těchto úseček 
musí protnout ve společném bodě těchto tří rovin. Bod $X$ tedy leží ve stejné 
rovině jako bod $P$ a mohu je spojit pro získání další čáry řezu.
\end{itemize}

\section*{Odchylky v prostoru}
Vždy hledáme odchylku $\phi$ v intervalu $\langle 0, 90^\circ \rangle$

\subsubsection*{Odchylka přímek}
U přímek různoběžných je potřeba najít rovinu ve které obě leží, poté už můžeme dopočítat

U různoběžných přímek nejprve najdeme rovinu se kterou jsou rovnoběžné (nejjednodušeji se najde tam, kde se přímky po posunutí protínají) a pak je obě do té roviny posunout jako je to uděláno na obrázku, kde je posunuta přímka $q$ do $q'$.

\begin{tikzpicture}[scale=1.4, line join=round, line cap=round]
    % Definice vrcholů pomocí 3D souřadnic (x, y, z)
    \coordinate (A) at (-1.5, 0,  1.5);
    \coordinate (B) at ( 1.5, 0,  1.5);
    \coordinate (C) at ( 1.5, 0, -1.5);
    \coordinate (D) at (-1.5, 0, -1.5);
    \coordinate (S) at ( 0,   0,  0);   % Střed podstavy
    \coordinate (V) at ( 0,   3,  0);   % Vrchol jehlanu
    \coordinate (Q1) at (-1.5, 0, 2.5);
    \coordinate (Q2) at (-1.5, 0, -2.5);
    \coordinate (P1) at (2.25, -1.5, 2.25);
    \coordinate (P2) at (-0.75, 4.5, -0.75);

    \coordinate (Q3) at (1.5, 0, 2.5);
    \coordinate (Q4) at (1.5, 0, -2.5);

    % Skryté hrany podstavy a zadní boční hrana (čárkovaně)
    \draw[dashed, thick] (A) -- (D) -- (C);
    \draw[dashed, thick] (D) -- (V);
    
    % Úhlopříčky podstavy a výška jehlanu (volitelné, pro názornost)
    \draw[dashed, thin, gray] (A) -- (C);
    \draw[dashed, thin, gray] (B) -- (D);
    \draw[dashed, thick, red] (S) -- (V) node[midway, right] {$v$};

    % Viditelné hrany (plnou čarou)
    \draw[thick] (A) -- (B) node[midway, below] {$a$};
    \draw[thick] (B) -- (C) node[midway, below right] {$a$};
    \draw[thick] (A) -- (V);
    \draw[thick] (B) -- (V);
    \draw[thick] (C) -- (V);

    %přímky
    \draw[thin] (Q1) -- (Q2);
    \draw[thin] (P1) -- (P2);
    \draw[thin, blue] (Q3) -- (Q4);

    % Vykreslení bodů a popisků
    \fill (A) circle (1pt) node[above left]  {$A$};
    \fill (B) circle (1pt) node[below right] {$B$};
    \fill (C) circle (1pt) node[right]       {$C$};
    \fill (D) circle (1pt) node[above left]  {$D$};
    \fill (S) circle (1pt) node[below]       {$S$};
    \fill (V) circle (1.5pt) node[above]     {$V$};
    
    \node[below left] at (Q1) {$q$};
    \node[above right] at (P1) {$p$};
    \node[below left] at (Q3) {$q'$};
\end{tikzpicture}

\subsubsection*{Odchylka přímky a roviny}
Definice: úhel mezi přímkou $p$ a jejím kolmým průmětem do této roviny (kolmý průmět je jako kdybychom se na to dívali zvrchu).

\subsubsection*{Odchylka rovin}
Odchylka dvou různoběžných rovin je úhel, který sírají jejich průsečnice s třetí rovinou, jež je kolmá k rovinám, tedy i k průsečnici.

\section*{Kolmost}
Přímka je kolmá k rovině, pokud je kolmá ke všem přímkám v dané rovině, které procházejí jejím průsečíkem. Pro ověření kolmosti stačí dokázat, že je přímka kolmá ke dvěma různobžkám ležícím v této rovině (u krychle prostě řeku, že je kolmá k tomu a kolmá k tomu, je tudíž kolmá  nějaké rovině)

Rovina $\rho$ je kolmá k rovině $\sigma$, pokud $\rho$ obsahuje přímku $p$, která je kolmá k rovině $\sigma$ nebo obráceně.

\section*{Vzdálenosti (metrické úlohy)}
\subsubsection*{Vzdálenost dvou přímek}
\begin{itemize}
    \item Různoběžné: vzdálenost je 0
    \item Rovnoběžné: zvolíme libovolný bod na přímce a spočítáme jeho vzdálenost od druhé přímky
    \item Mimoběžné: počítáme jejich nejkratší vzdálenost, která se nazývá osa
\end{itemize}


\end{document}
